\documentclass{article}
\usepackage[a4paper]{geometry}
\usepackage{fancyhdr}
\pagestyle{fancy}
\lhead{Handelsregime}
\rhead{März 2026}
\begin{document}
\section{Handelsregime}
\emph{Handelsregime} beschreiben Organisationen, welche versuchen den Welthandel zu regulieren.
 
\subsection{Die WTO} 
Die größte davon ist die \emph{WTO}, die \emph{World Trade Organisation} beziehungsweise die \emph{Welthandelsorganisation}; sie umfasst ungefähr $98\%$ des Welthandels. Auf unterschiedliche Art und Weisen versucht sie dafür zu sorgen, dass verbindliche Regeln geschaffen werden, unter welchen kein Mitgliedsstaat von anderen diskriminert wird und der Handel verbessert wird, z.\,B. indem Zölle gesenkt werden.
 
Im \emph{DSB} der WTO, dem \emph{Dispute Settlement Body}, überprüfen Expertengremien Streitigkeiten zwischen Staaten. Das Urteil sollte formal bindend sein.
 
\subsection{Grenzen} 
Die WTO hat aber nahezu keine Macht, ihre Entscheidungen umzusetzen, dies müssen die Mitgliedsstaaten (mit der benötigeten Macht) tun. So können mächtige Staaten auch Urteile und Regeln ignorieren. Dies tun die USA und China, auch in letzter Zeit, gerne.
 
Zudem verfolgt die WTO ein Konsensprinzip, welches bei $> 150$ Mitgliedsstaaten lange dauert und vermehrt zu Handlungsunfähigkeit führt. So passiert es seit 2001 in der \emph{Doha-Runde}.
\end{document}